\documentclass[11pt,a4paper]{article}
\usepackage[utf8]{inputenc}
\usepackage[T1]{fontenc}
\usepackage[portuguese]{babel}
\usepackage{lmodern}
\usepackage[margin=2.3cm]{geometry}
\usepackage{tabularx,longtable,booktabs,array,enumitem}
\usepackage[table]{xcolor}
\usepackage{tcolorbox}
\usepackage{hyperref}
\definecolor{roamly}{HTML}{0E6E68}
\definecolor{alerta}{HTML}{B3261E}
\definecolor{formal}{HTML}{2E7D32}
\hypersetup{colorlinks=true,linkcolor=roamly,urlcolor=roamly}
\newtcolorbox{aviso}[1][]{colback=alerta!6,colframe=alerta,arc=2pt,boxrule=0.8pt,fonttitle=\bfseries,title=#1}
\newcolumntype{L}{>{\raggedright\arraybackslash}X}
\renewcommand{\arraystretch}{1.3}
\setlist{nosep,leftmargin=*}
\newcommand{\bolt}{\textcolor{alerta}{[BD]}}
\newcommand{\ent}[1]{\textit{Entrevistador:} #1}

\title{\textbf{Roamly}\\[2pt]\large Entrevista a agente de viagens: A0\\\normalsize Registo da entrevista (caso fictício)}
\author{}\date{}
\begin{document}
\maketitle

\begin{aviso}[Caso 1: Profissional de uma agência de viagens] 
Tema: \textbf{preocupações com a IA e impacto no negócio}. 
Anónimo: sem nomes de pessoas, empresas, plataformas ou 
cidades. O código \texttt{A0} é reservado a este caso.
\end{aviso}

%==========================================================
\section{Perfil e situação da entrevista}
\begin{tabularx}{\textwidth}{@{}p{4.3cm}L@{}}
\toprule
Código & A0 \\
Papel & Agente de viagens de balcão \\
Perfil & 12 anos de experiência na mesma agência de viagens. \\
Tipo de tarefa & Normal, no posto de trabalho: manhã de trabalho \textbf{observada} no balcão, 
mais \textbf{relatos retrospetivos} de casos concretos \\
Local e duração & Na agência, 11\,h15--12\,h; \(\approx\)35\,min de observação ativa \\
Registo & Transcrição do conteúdo considerado necessário (com consentimento); sem fotos do balcão \\
\bottomrule
\end{tabularx}

%==========================================================
\section{Siglas}
{\small
\begin{tabularx}{\textwidth}{@{}p{2.6cm}L@{}}
\toprule
\textbf{Sigla} & \textbf{Significado} \\ \midrule
\multicolumn{2}{@{}l}{\textit{Identificadores}} \\
A0 & Código do agente entrevistado (caso fictício) \\
AE01--AE08 & Registos da entrevista ao agente (observações e relatos) \\
A0-01--A0-15 & Notas da sessão de interpretação do caso A0 \\ \midrule
\multicolumn{2}{@{}l}{\textit{Tipos de registo da entrevista}} \\
O & Observado pelo entrevistador \\
R & Relato de caso concreto contado durante a observação \\
F/OP & Fala geral ou opinião, sem caso concreto (fora dos modelos) \\ \midrule
\multicolumn{2}{@{}l}{\textit{Tipos de nota de interpretação}} \\
OBS & Observação \\
INS & \textit{Insight} \\
CUL & Influência cultural \\
Q & Pergunta por esclarecer \\
BD & \textit{Breakdown}: falha ou atrito no trabalho \\
OP & Opinião geral \\ \midrule
\multicolumn{2}{@{}l}{\textit{Outras}} \\
\bolt{} & Marca de \textit{breakdown} nos registos AE \\
IA & Inteligência artificial \\
\bottomrule
\end{tabularx}}

%==========================================================
\section{Conversa convencional (\(\approx\)10\,min)}
Tipo F/OP: contextualiza mas \textbf{não entra nos modelos} sem caso concreto.

\begin{tabularx}{\textwidth}{@{}p{1.6cm}L@{}}
\toprule
\textit{Entrev.} & «Quais as principais funções que desempenha?» \\
A0 & «Aconselho, reservo, e resolvo quando corre mal.» \\
\textit{Entrev.} & «Como vê a IA no seu trabalho?» \\
A0 & «\textcolor{formal}{Por um lado, tenho medo de que os clientes deixem de precisar de nós, por outro
o que vejo é que deixa as pessoas ainda mais confusas.}» (\textbf{OP}) \\
\textit{Entrev.} & «O que mudou no negócio nos últimos anos?» \\
A0 & «\textcolor{formal}{Os clientes chegam informados e com ideias definidas, 
a maioria mal informados e sem tentarem entender se de facto é possível fazer mesmo aquilo.
Acabamos por parecer «desmancha-prazeres» por tentar explicar que a logística não tem grande sentido.}» (\textbf{OP}) \\
\bottomrule
\end{tabularx}

\paragraph{Transição.} «Vou ficar aqui a seguir o seu trabalho esta manhã e a perguntar o que não perceber. Se aparecer um cliente, continue como faz normalmente.»

%==========================================================
\section{Entrevista contextual: o que foi observado e contado}
Observação ativa no balcão: 11\,h25--12\,h. \textbf{AE} = registo da entrevista ao agente; \textbf{O} = observado; \textbf{R} = relato de caso concreto contado durante a observação; \textbf{F/OP} = geral (fora dos modelos).

{\small
\begin{longtable}{@{}p{1.1cm}p{0.7cm}p{13.4cm}@{}}
\toprule
\textbf{ID} & \textbf{Tipo} & \textbf{Registo} \\ \midrule \endhead
AE01 & O & 11\,h25. Abre o sistema com algumas conversas pendentes e a caixa de correio 
partilhada. Um papel autocolante no monitor 
com «ligar a cliente X». \\

AE02 & O & Lê um pedido novo no correio: «vi opções de alojamento parecidas,
por menos 10 ou 20 euros por noite». Suspira: «É isto, lá vamos nós.» \bolt{} 
pressão sobre o preço. \\

AE03 & O & Entra um casal, com um roteiro impresso: algumas ideias, voos
e sítios onde gostariam de ficar. Querem «algo deste género». 
A0 \emph{pergunta o orçamento, as datas flexíveis e «o que não abdicam»}. \\

AE04 & O & Ao ler o roteiro, A0 mostra que duas cidades ficam a 6\,h de comboio e 
que um monumento fecha no dia previsto. O casal troca olhares desagradado. \\


AE05 & O & Começa a proposta: copia dados do portal para o modelo da agência 
(documento com logótipo); diz que a envia por email «o mais rápido possível». \\

AE06 & O & O casal agradece e diz que vai «pensar e comparar online». \\

AE07 & R & \textit{Caso:} um cliente recusou pagar a taxa de serviço
«porque o site não cobra». A0 acabou por perder o cliente. \bolt{} \\

AE08 & O & Mostra o caderno de fornecedores: «hotel X: não fiar na foto; pedir 
quarto nos pisos de cima». «Isto não está em lado nenhum.» \\

\bottomrule
\end{longtable}}

%==========================================================
\section{Fecho: resumo validado com A0}
\begin{itemize}
  \item «O seu valor é aconselhar, responsabilizar-se e resolver quando corre mal.» 
  A0 acrescenta: «e conhecer os fornecedores.»
  \item «Gasta muito tempo a procurar e copiar dados entre sistemas.» A0: 
  «sim, grande parte da manhã.»
  \item «O ChatGPT preocupa-a quando parece certa e não é.» A0 corrige: «e quando o 
  supervisor pensa que ele (ChatGPT) me substitui, mas não gosta quando o relatório
  sai confuso.»
\end{itemize}

%==========================================================
\section{Notas da sessão de interpretação}
\begin{aviso}
  Notas sobre o software usado foram assumidas e outras geradas pela nossa imaginação
  uma vez que nos pediram para não divulgar o seu funcionamento.
\end{aviso}
Tipos: \textbf{OBS} observação; \textbf{INS} insight; \textbf{CUL} influência cultural; \textbf{Q} pergunta por esclarecer; \textbf{BD} breakdown; \textbf{OP} opinião geral.

{\small
\begin{longtable}{@{}p{1.5cm}p{1.1cm}p{11.8cm}@{}}
\toprule
\textbf{ID} & \textbf{Tipo} & \textbf{Nota} \\ \midrule \endhead
A0-01 & OBS & Trabalha com vários sítios em paralelo: sistema com conversas pendentes, correio partilhado e papel autocolante no monitor (AE01). \\
A0-02 & CUL & Os clientes comparam preços online e usam isso como pressão (AE02, AE07). \\
A0-03 & BD & Cliente recusa a taxa de serviço porque «o site não cobra»; A0 perde o cliente (AE07). \\
A0-04 & OBS & Os clientes chegam com roteiros impressos e ideias definidas (AE03). \\
A0-05 & BD & O roteiro do cliente tinha logística inviável; A0 tem de o desfazer e o casal fica desagradado. Liga-se à opinião «desmancha-prazeres» da conversa (AE04). \\
A0-06 & INS & A0 faz perguntas de descoberta («o que não abdicam?») a que o roteiro do cliente não respondia (AE03). \\
A0-07 & OBS & A proposta é copiada à mão do portal para o modelo da agência (AE05). \\
A0-08 & INS & Os clientes usam a agência para se aconselhar e depois comparam sozinhos online (AE06). \\
A0-09 & CUL & O supervisor vê o ChatGPT como substituto de A0; A0 vê-o como um risco que lhe cai em cima (fecho). \\
A0-10 & OBS & Conhecimento tácito sobre fornecedores guardado num caderno pessoal (AE08). \\
A0-11 & CUL & Identidade: «aconselho, reservo, resolvo» e «conhecer os fornecedores»; teme ser reduzido(a) a intermediário (conversa, fecho). \\
A0-12 & OP & «Tenho medo de que os clientes deixem de precisar de nós» (sem caso concreto; fica como contexto cultural). \\
A0-13 & OP & «A IA deixa as pessoas ainda mais confusas» (sem caso concreto observado; fica como contexto). \\
A0-14 & Q & Como saberia o cliente que a proposta foi revista por uma pessoa? Perguntar a mais agentes. \\
A0-15 & Q & Quanto da receita vem de comissões versus taxas de serviço? Pedir em entrevista futura (não observado). \\
\bottomrule
\end{longtable}}

\end{document}
